% ==STYLE-BEGIN==
\documentclass[UTF8,11pt]{ctexart}
\usepackage{amsmath,amssymb,amsthm}
\usepackage[a4paper,margin=2.4cm,top=2.3cm,bottom=2.3cm]{geometry}
\usepackage{fancyhdr}
\usepackage{enumitem}
\usepackage{hyperref}
\hypersetup{hidelinks}

% ---- 版式 ----
\linespread{1.15}
\setlength{\parskip}{0.35em}
\setlength{\parindent}{1.6em}
\setlist{nosep,leftmargin=2em}

% ---- 页眉 ----
\pagestyle{fancy}
\fancyhf{}
\fancyhead[L]{\small\coursename}
\fancyhead[R]{\small\thepage}
\renewcommand{\headrulewidth}{0.4pt}

% ---- 定理环境：中文用正体，标题加粗，按小节共用一条流水号 ----
\newtheoremstyle{cnthm}{0.7em}{0.7em}{\normalfont}{}{\bfseries}{}{0.5em}{}
\theoremstyle{cnthm}
\newtheorem{thm}{定理}[section]
\newtheorem{lem}[thm]{引理}
\newtheorem{prop}[thm]{命题}
\newtheorem{cor}[thm]{推论}
\newtheorem{defn}[thm]{定义}
\newtheorem{rem}[thm]{注}
\newtheorem{exa}[thm]{例}
\newtheorem{conj}[thm]{猜想}
\renewcommand{\proofname}{\bfseries 证明}
\providecommand{\coursename}{}
% ==STYLE-END==
\renewcommand{\coursename}{数学分析\,·\,1029}
\title{数学分析 1029}
\date{}
\begin{document}
\maketitle

\section{高阶导数与泰勒级数}

\subsection{高阶导数}

$f$ 在 $x_0$ 处可导 $\Rightarrow$ $f$ 在 $x_0$ 处连续，但推不出 $f$ 在 $x_0$ 附近处处可导。

若 $f$ 在 $x_0$ 的某个开邻域 $U$ 内点点可导，则 $\forall x\in U$，$f'(x)$ 存在，于是得到 $U$ 上的函数
\[
f':\ U\to\mathbb{R},\qquad x\longmapsto f'(x).
\]
若 $f'$ 在 $x_0$ 处可导，则记 $(f'(x))'|_{x=x_0}$ 为 $f''(x_0)$ 或 $f^{(2)}(x_0)$。

\begin{rem}
（1）$f$ 在 $x_0$ 处可导 $\Rightarrow$ $f$ 在 $x_0$ 处连续；

（2）二阶可导 $\Rightarrow$ $x_0$ 的某个开区间内点点可导；

（3）三阶可导 $\Rightarrow$ 二阶可导、一阶可导；$\cdots$；$k$ 阶可导 $\Rightarrow$ $k-1,k-2,\cdots,1$ 阶都可导。
\end{rem}

若 $f,g$ 在 $x_0$ 处 $n$ 阶可导，则
\[
(f\pm g)^{(n)}(x_0)=f^{(n)}(x_0)\pm g^{(n)}(x_0).
\]

\begin{rem}
\[
f^{(n)}(x_0)=\lim_{x\to x_0}\frac{f^{(n-1)}(x)-f^{(n-1)}(x_0)}{x-x_0}
=\left.(f^{(n-1)})'\right|_{x=x_0},
\]
更一般地 $f^{(n)}(x_0)=\left.(f^{(k)})^{(n-k)}\right|_{x=x_0}$（$0\leq k\leq n$）。
\end{rem}

\subsection{乘积的高阶导数（Leibniz 公式）}

\begin{thm}[Leibniz]
若 $f,g$ 在 $x_0$ 处 $n$ 阶可导，则
\[
(fg)^{(n)}=\sum_{k=0}^{n}\binom{n}{k}f^{(n-k)}g^{(k)} .
\]
\end{thm}

\begin{proof}
对 $n$ 归纳。$n=1$ 即 $(fg)'=f'g+fg'$。设对 $n$ 成立，则
\[
(fg)^{(n+1)}=\Bigl(\sum_{k=0}^{n}\binom{n}{k}f^{(n-k)}g^{(k)}\Bigr)'
=\sum_{k=0}^{n}\binom{n}{k}f^{(n+1-k)}g^{(k)}+\sum_{k=0}^{n}\binom{n}{k}f^{(n-k)}g^{(k+1)},
\]
把第二个和式中的下标换成 $k+1$ 并利用 $\binom{n}{k}+\binom{n}{k-1}=\binom{n+1}{k}$，即得 $n+1$ 的情形。
\end{proof}

\subsection{用算子再看这个公式}

引入算子 $\partial$：把函数映成它的导函数，
\[
\partial:\ f\longmapsto f',\qquad\text{于是}\quad f^{(n)}=\partial^{n}f .
\]
\[
\partial(fg)=\partial(f)\,g+f\,\partial(g)
=(\partial\otimes1)(fg)+(1\otimes\partial)(fg)
=(\partial\otimes1+1\otimes\partial)(fg),
\]
\[
\partial^{n}(fg)=(\partial\otimes1+1\otimes\partial)^{n}(fg).
\]
类比 $(a+b)^{2}=a^{2}+ab+ba+b^{2}$：只有 $ab=ba$ 时才有 $(a+b)^{2}=a^{2}+2ab+b^{2}$。这里要说明
\[
(\partial\otimes1)(1\otimes\partial)=(1\otimes\partial)(\partial\otimes1),
\]
两者作用的结果都是 $f'g'$：
\[
(\partial\otimes1)(1\otimes\partial)(fg)=(1\otimes\partial)(\partial\otimes1)(fg)=f'g'.
\]
于是可以按二项式展开：
\[
(\partial\otimes1+1\otimes\partial)^{n}(fg)
=\sum_{k=0}^{n}\binom{n}{k}(\partial\otimes1)^{k}(1\otimes\partial)^{n-k}(fg)
=\sum_{k=0}^{n}\binom{n}{k}f^{(k)}g^{(n-k)} .
\]

\subsection{两个例子}

\begin{exa}
$y=\sin x$：
\[
y'=\cos x=\sin\Bigl(x+\frac{\pi}{2}\Bigr),\quad
y''=\sin\Bigl(x+2\cdot\frac{\pi}{2}\Bigr),\quad\cdots,\quad
y^{(n)}=\sin\Bigl(x+n\frac{\pi}{2}\Bigr).
\]
\end{exa}

\begin{exa}
$y=x^{n}$：
\[
y'=nx^{n-1},\quad y''=n(n-1)x^{n-2},\quad\cdots,\quad
y^{(k)}=n(n-1)\cdots(n-k+1)x^{n-k},\quad y^{(n)}=n! .
\]
\end{exa}

\subsection{从多项式到泰勒系数}

设
\[
f(x)=a_nx^{n}+a_{n-1}x^{n-1}+\cdots+a_1x+a_0 ,
\]
已知 $\forall k$ 的 $f^{(k)}(0)$，并约定 $f^{(0)}(x)=f(x)$。先看二次的情形：
\[
f(x)=a_2x^{2}+a_1x+a_0\ \Longrightarrow\ a_0=f(0);
\]
\[
f'(x)=2a_2x+a_1\ \Longrightarrow\ a_1=f'(0);\qquad
f''(x)=2a_2\ \Longrightarrow\ a_2=\frac12f''(0).
\]
推到一般情形求 $a_m$：$f^{(m)}(x)$ 中只有 $a_mx^{m}$ 求 $m$ 次导后不留 $x$，其余各项都含 $x$ 的幂，于是令 $x=0$ 得 $a_m=\dfrac{1}{m!}f^{(m)}(0)$。更一般的写法：把 $f^{(k)}(x)$ 写成
\[
f^{(k)}(x)=k!\,a_k+x\,h(x),
\]
令 $x=0$ 即得 $a_k=\dfrac{f^{(k)}(0)}{k!}$。因此
\[
f(x)=\sum_{k=0}^{n}\frac{f^{(k)}(0)}{k!}x^{k}.
\]
（求导用来消去低次项，令 $x=0$ 用来消去高次项。）

\subsection{一般函数能这么做吗}

如果 $g(x)$ 可在 $0$ 点求任意阶导数，令
\[
b_k=\frac{g^{(k)}(0)}{k!},
\]
构造多项式 $b_0+b_1x+\cdots+b_kx^{k}+\cdots+b_nx^{n}$。想用它逼近 $g(x)$，只需要验证
\[
g(x)-\bigl(b_0+b_1x+\cdots+b_nx^{n}\bigr)
\]
与 $x^{n}$ 比较阶数——这就是 Taylor 级数要处理的问题。

\begin{exa}
$f(x)=x^{2}\sin x$，求 $f^{(n)}(0)$。由 Leibniz 公式只有 $k=0,1,2$ 三项非零，而在 $x=0$ 处前两项为零，故
\[
f^{(n)}(0)=\binom{n}{2}\cdot 2\cdot\sin\frac{(n-2)\pi}{2}
=n(n-1)\sin\frac{(n-2)\pi}{2}.
\]
\end{exa}

\begin{exa}
$y=\dfrac{2x+3}{(x+1)(x+2)}$，求 $y^{(n)}(x)$。先把分式拆成部分分式
\[
y=\frac{1}{x+1}+\frac{1}{x+2},
\]
再用 $\bigl((x+c)^{-1}\bigr)^{(n)}=(-1)^{n}n!\,(x+c)^{-(n+1)}$ 逐项求导。
\end{exa}

\begin{rem}
求高阶导时，如果能先把乘除换成加减（部分分式、取对数等），会省事很多。
\end{rem}

\section{隐函数求导}

\begin{defn}[隐函数]
设 $F(x,y)=0$。如果通过限制取值范围的方式，使得范围内 $\forall x$ 都有唯一的 $y$ 与之对应，即 $F\bigl(x,y(x)\bigr)=0$，则称函数 $x\mapsto y(x)$ 为由 $F(x,y)=0$ 决定的隐函数。
\end{defn}

在 $F\bigl(x,y(x)\bigr)=0$ 两边对 $x$ 求导，得
\[
\frac{\partial F}{\partial x}+\frac{\partial F}{\partial y}\,y'(x)=0,
\]
只要 $\dfrac{\partial F}{\partial y}\neq0$，就可解出 $y'(x)$。

\begin{exa}
设 $y=y(x)$ 由 $e^{y}+xy=e$ 决定，求 $y'(x)|_{x=0}$。

两边同时对 $x$ 求导：$e^{y}y'+y+xy'=0$，故
\[
y'\Big|_{x=0}=-\left.\frac{y}{e^{y}+x}\right|_{x=0}=-\frac{1}{e}
\qquad(\text{当 }x=0\text{ 时 }y=1).
\]
\end{exa}

\begin{exa}
$f(x)=\arctan x$，求 $f^{(n)}(0)$。令 $y=\dfrac{1}{1+x^{2}}$，则 $y^{(n)}(0)=f^{(n+1)}(0)$，问题化为求 $y$ 的高阶导。

由 $y(1+x^{2})=1$，两边同时求 $n$ 阶导（用 Leibniz 公式）得
\[
y^{(n)}\bigl(1+x^{2}\bigr)+2nx\,y^{(n-1)}+2\binom{n}{2}y^{(n-2)}=0 .
\]
令 $x=0$：
\[
y^{(n)}(0)+2\binom{n}{2}y^{(n-2)}(0)=0,
\]
即 $y^{(n)}(0)=-n(n-1)y^{(n-2)}(0)$。由 $y(0)=1$、$y'(0)=0$ 可依次推出所有 $y^{(n)}(0)$，从而得到 $f^{(n)}(0)$。
\end{exa}

\section{参数方程求导}

设
\[
\begin{cases}
x=\varphi(t),\\
y=\psi(t).
\end{cases}
\]
若在某开区间内 $\varphi(t)$ 严格单调且 $\varphi'(t)\neq0$，则 $x$ 可以反解出 $t=g(x)$，代入 $y=\psi(t)$ 得 $y=\psi\bigl(g(x)\bigr)$。

此时 $y$ 是 $x$ 的函数，求 $y'(x)$：
\[
y'(x)=\psi'\bigl(g(x)\bigr)\cdot g'(x)=\psi'(t)\cdot g'(x),
\]
其中由反函数求导可知 $g'(x)=\dfrac{1}{\varphi'(t)}$。于是
\[
y'(x)=\left.\frac{\psi'(t)}{\varphi'(t)}\right|_{t=g(x)} .
\]

\section{微分}

\begin{thm}[可导与可微]
$f$ 在 $x_0$ 处可导
\[
\iff f(x)=f(x_0)+A\Delta x+o(\Delta x)
\iff f(x)-f(x_0)=A\Delta x+o(\Delta x)=\Delta f .
\]
\end{thm}

忽略高阶无穷小 $o(\Delta x)$，记
\[
\mathrm{d}f=A\Delta x=A\,\mathrm{d}x
\qquad\Longrightarrow\qquad
\frac{\mathrm{d}f}{\mathrm{d}x}=f' .
\]
对参数方程 $\begin{cases}y=\psi(t),\\ x=\varphi(t)\end{cases}$，同样有
\[
\frac{\mathrm{d}y}{\mathrm{d}x}=\frac{\mathrm{d}y/\mathrm{d}t}{\mathrm{d}x/\mathrm{d}t}
=\frac{\psi'(t)}{\varphi'(t)} .
\]

\begin{defn}[微分]
$\mathrm{d}f$ 称为 $f$ 的微分。若 $f$ 在 $x_0$ 附近能写成
\[
f(x)=f(x_0)+A\Delta x+o(\Delta x),
\]
则称 $f$ 在 $x_0$ 处可微。可微 $\iff$ 可导。
\end{defn}

\end{document}
